\documentclass[aspectratio=169, 10pt]{beamer}
\usepackage{beamer_theme_cienciadados2026}

\title[Aula 05 - Fundamentos de Python]{Ciência dos Dados I: Análise Exploratória}
\subtitle{Aula 05: Fundamentos de Python: Sintaxe, Strings e Formatação Moderna}
\author{Prof. Dr. Ney Lemke}
\institute{UNESP -- Instituto de Biociências de Botucatu}
\date{Versão 2026}

\begin{document}

\frame{\titlepage}

\begin{frame}{Sumário da Aula}
  \tableofcontents
\end{frame}

\section{Filosofia e Legibilidade}

\begin{frame}{O Zen do Python (\texttt{import this})}
  \begin{itemize}
    \item \textit{"Bonito é melhor que feio."}
    \item \textit{"Explícito é melhor que implícito."}
    \item \textit{"Simples é melhor que complexo."}
    \item \textit{"Legibilidade conta."}
  \end{itemize}
  \begin{block}{PEP 8: Guia de Estilo Oficial}
    Identação obrigatória de 4 espaços (nunca misturar com tabs), nomenclatura \texttt{snake\_case} para funções e variáveis, e espaçamento consistente ao redor de operadores.
  \end{block}
\end{frame}

\section{Manipulação e Sanitização de Strings}

\begin{frame}[fragile]{Métodos Essenciais para Limpeza de Dados Textuais}
  \begin{itemize}
    \item \texttt{str.strip()}: Remove espaços em branco e quebras de linha nas extremidades.
    \item \texttt{str.split(';')}: Divide o texto em uma lista de strings baseando-se em um delimitador.
    \item \texttt{str.join(lista)}: Concatena uma lista de palavras usando um separador.
    \item \texttt{str.replace(antigo, novo)}: Substitui ocorrências de padrões no texto.
  \end{itemize}
  \begin{lstlisting}
registro = "   P01 ; Glicose ; 115.4 mg/dL \n"
campos = [c.strip() for c in registro.split(';')]
# Resultado: ['P01', 'Glicose', '115.4 mg/dL']
  \end{lstlisting}
\end{frame}

\section{F-Strings e Fatiamento de Sequências}

\begin{frame}[fragile]{Interpolação Moderna e Slicing}
  \begin{columns}
    \begin{column}{0.5\textwidth}
      \textbf{F-strings Avançadas}
      \begin{lstlisting}
valor = 1420.5
taxa = 0.0825
print(f"R$ {valor:10.2f}")
print(f"Taxa: {taxa:.1%}")
      \end{lstlisting}
    \end{column}
    \begin{column}{0.48\textwidth}
      \textbf{Fatiamento: \texttt{[start:stop:step]}}
      \begin{lstlisting}
seq = "DNA_TTAGCA"
print(seq[4:])    # 'TTAGCA'
print(seq[::-1])  # 'ACGATT_AND'
      \end{lstlisting}
    \end{column}
  \end{columns}
\end{frame}

\begin{frame}{Resumo da Aula}
  \begin{itemize}
    \item Strings são sequências imutáveis fundamentais para identificadores, categorias e logs.
    \item F-strings oferecem precisão decimal e formatação clara para relatórios analíticos.
    \item \textbf{Prática}: Higienização de cadastros e extração de padrões em dados brutos.
  \end{itemize}
\end{frame}

\end{document}
